\documentclass[11pt,a4paper]{article}
\usepackage[utf8]{inputenc}
\usepackage[T1]{fontenc}
\usepackage{geometry}
\usepackage{amsmath}
\usepackage{amssymb}
\usepackage{graphicx}
\usepackage{xcolor}
\usepackage{fancyhdr}
\usepackage{titlesec}
\usepackage{listings}
\usepackage{tcolorbox}
\usepackage{enumitem}
\usepackage{booktabs}
\usepackage{array}
\usepackage{longtable}
\usepackage{hyperref}
\usepackage{tikz}
\usetikzlibrary{arrows.meta,decorations.markings,calc,angles,quotes}

% Page geometry
\geometry{margin=2.5cm}

% Colors
\definecolor{titlecolor}{RGB}{44,62,80}
\definecolor{sectioncolor}{RGB}{52,73,94}
\definecolor{codebackground}{RGB}{244,244,244}
\definecolor{codeborder}{RGB}{52,152,219}
\definecolor{emphcolor}{RGB}{231,76,60}

% Header and footer
\pagestyle{fancy}
\fancyhf{}
\fancyhead[L]{\small Space Debris Removal - Background Tutorial}
\fancyhead[R]{\small IIT Kanpur Hackathon 2025}
\fancyfoot[C]{\thepage}
\renewcommand{\headrulewidth}{0.4pt}

% Title formatting
\titleformat{\section}
{\color{titlecolor}\Large\bfseries}
{\thesection}{1em}{}[\vspace{-0.5em}\color{codeborder}\titlerule]

\titleformat{\subsection}
{\color{sectioncolor}\large\bfseries}
{\thesubsection}{1em}{}

\titleformat{\subsubsection}
{\color{sectioncolor}\normalsize\bfseries}
{\thesubsubsection}{1em}{}

% Code listing style
\lstset{
    basicstyle=\ttfamily\small,
    backgroundcolor=\color{codebackground},
    frame=single,
    framerule=1pt,
    rulecolor=\color{codeborder},
    breaklines=true,
    columns=fullflexible,
    keepspaces=true,
    xleftmargin=10pt,
    xrightmargin=10pt,
    framexleftmargin=5pt,
    framexrightmargin=5pt,
}

% Hyperlink setup
\hypersetup{
    colorlinks=true,
    linkcolor=blue,
    urlcolor=blue,
    citecolor=blue
}

\begin{document}

% Title page
\begin{titlepage}
    \centering
    \vspace*{2cm}

    {\Huge\bfseries\color{titlecolor} Space Debris Removal\par}
    \vspace{0.5cm}
    {\LARGE\color{sectioncolor} Background for Robotics \& Controls Engineers\par}
    \vspace{2cm}

    {\Large IIT Kanpur Research Symposium 2025\par}
    \vspace{0.5cm}
    {\large 24-Hour Hackathon\par}

    \vfill

    {\large Theme: Aerospace, Space Science \& Engineering\par}
    \vspace{0.5cm}
    {\large Mission Concept for Mitigating Space Debris in LEO\par}

    \vspace{1cm}
    {\large \today\par}
\end{titlepage}

\tableofcontents
\newpage

\section{Space Debris Problem Overview}

\subsection{What is Space Debris?}

\begin{itemize}[leftmargin=*]
    \item \textbf{Defunct satellites}: No longer operational spacecraft
    \item \textbf{Rocket bodies}: Upper stages left in orbit
    \item \textbf{Fragmentation debris}: Pieces from collisions and explosions
    \item \textbf{Mission-related objects}: Lens caps, separation mechanisms, etc.
\end{itemize}

\subsection{LEO Characteristics}

\begin{itemize}[leftmargin=*]
    \item \textbf{Altitude}: 200--2000 km above Earth
    \item \textbf{Orbital velocity}: $\sim$7.8 km/s (28,000 km/h!)
    \item \textbf{Orbital period}: $\sim$90--120 minutes
    \item \textbf{Atmospheric drag}: Still present, causes gradual decay
\end{itemize}

\subsection{Why It Matters}

\begin{itemize}[leftmargin=*]
    \item Over 34,000 tracked objects $>$10 cm
    \item $\sim$1 million objects between 1--10 cm
    \item Collision at orbital speeds can destroy spacecraft
    \item Even small debris (1 cm) can be catastrophic
\end{itemize}

\section{Orbital Mechanics Fundamentals}

\subsection{Kepler's Laws}

As a controls engineer, you need to understand that satellites follow predictable paths:

\subsubsection{Elliptical Orbits}
\begin{itemize}[leftmargin=*]
    \item Orbits are ellipses with Earth at one focus
    \item Described by semi-major axis ($a$) and eccentricity ($e$)
\end{itemize}

\subsubsection{Orbital Parameters (State Vector)}

\begin{lstlisting}
Position: [x, y, z]
Velocity: [vx, vy, vz]
\end{lstlisting}

Or using \textbf{Keplerian Elements}:
\begin{itemize}[leftmargin=*]
    \item \textbf{$a$}: Semi-major axis (orbit size)
    \item \textbf{$e$}: Eccentricity (orbit shape: 0 = circular, $<$1 = ellipse)
    \item \textbf{$i$}: Inclination (tilt relative to equator)
    \item \textbf{$\Omega$}: Right Ascension of Ascending Node (RAAN)
    \item \textbf{$\omega$}: Argument of Periapsis
    \item \textbf{$\nu$}: True Anomaly (position in orbit)
\end{itemize}

\begin{figure}[htbp]
\centering

% Define colors globally for this figure
\definecolor{orbitcolor}{RGB}{52,152,219}
\definecolor{earthcolor}{RGB}{46,125,50}
\definecolor{equatorcolor}{RGB}{255,152,0}
\definecolor{refcolor}{RGB}{158,158,158}

\begin{tikzpicture}[scale=1.1, >=Stealth]

    % Equatorial plane (reference plane)
    \fill[equatorcolor!10] (0,0) ellipse (4.2cm and 0.45cm);
    \draw[equatorcolor, very thick, dashed] (0,0) ellipse (4.2cm and 0.45cm);

    % Reference direction (Vernal Equinox)
    \draw[->, refcolor, very thick] (0,0) -- (4.5,0) node[right, font=\normalsize] {$\gamma$};

    % Earth at focus
    \shade[ball color=earthcolor] (0,0) circle (0.3cm);

    % Orbital ellipse (tilted for 3D effect)
    \begin{scope}[rotate=15, yshift=0.3cm]
        % Draw the orbit
        \draw[orbitcolor, line width=2pt] (0,0) ellipse (3.3cm and 2.1cm);

        % Periapsis point
        \filldraw[red!70!black] (-3.3,0) circle (2.5pt);

        % Apoapsis point
        \filldraw[red!70!black] (3.3,0) circle (2.5pt);

        % Semi-major axis (a)
        \draw[<->, blue!80!black, very thick] (-3.3,-0.6) -- (3.3,-0.6);
        \node[blue!80!black, font=\Large\bfseries] at (0,-0.9) {$a$};

        % Ascending node
        \filldraw[green!50!black] (2.7,-1.25) circle (3pt);

        % Descending node
        \filldraw[green!30!black] (-2.7,1.25) circle (2pt);

        % Line of nodes (dashed)
        \draw[green!40!black, dashed, thick] (-3.3,1.5) -- (3.3,-1.5);

        % Satellite position
        \filldraw[orange!80!black] (1.7,1.4) circle (3.5pt);

        % True anomaly (ν) - angle from periapsis to satellite
        \draw[->, violet!80!black, very thick] (0,0) -- (1.7,1.4);
        \draw[violet!80!black, very thick, ->] (1.1,0) arc (0:40:1.1);
        \node[violet!80!black, font=\Large\bfseries] at (1.4,0.45) {$\nu$};

        % Argument of periapsis (ω)
        \draw[brown!80!black, very thick, ->] (1.4,-0.65) arc (-25:175:0.75);
        \node[brown!80!black, font=\Large\bfseries] at (-0.6,-0.95) {$\omega$};
    \end{scope}

    % Inclination angle (i)
    \draw[red!80!black, very thick, ->] (2.8,0.1) arc (5:25:2.8);
    \node[red!80!black, font=\Large\bfseries] at (3.3,0.7) {$i$};

    % RAAN (Ω)
    \draw[green!50!black, very thick, ->] (3.2,0) arc (0:-32:3.2);
    \node[green!50!black, font=\Large\bfseries] at (3.7,-1.0) {$\Omega$};

    % North direction indicator
    \draw[->, thick, gray!80!black] (-4.5,2.2) -- (-4.5,3.0);
    \node[gray!80!black, left, font=\small] at (-4.5,2.6) {$z$};

\end{tikzpicture}

\vspace{0.3cm}

% Legend below the figure
\begin{tcolorbox}[colback=blue!5!white,colframe=codeborder,width=0.95\textwidth,arc=2mm]
\small
\textbf{Legend:}
\begin{itemize}[leftmargin=*, nosep]
    \item \textcolor{earthcolor}{\textbf{Green sphere}}: Earth (at focus of ellipse)
    \item \textcolor{equatorcolor}{\textbf{Orange dashed ellipse}}: Equatorial plane (reference plane)
    \item \textcolor{orbitcolor}{\textbf{Blue ellipse}}: Orbital plane (satellite's path)
    \item \textcolor{refcolor}{$\gamma$}: Vernal equinox direction (reference for $\Omega$)
    \item \textcolor{orange!80!black}{\textbf{Orange dot}}: Satellite position
    \item \textcolor{red!70!black}{\textbf{Red dots}}: Periapsis (closest) and apoapsis (farthest) points
    \item \textcolor{green!50!black}{\textbf{Green dots}}: Ascending node (orbit crosses equatorial plane going north)
\end{itemize}

\textbf{Six Keplerian Elements:}
\begin{itemize}[leftmargin=*, nosep]
    \item \textcolor{blue!80!black}{$a$}: \textbf{Semi-major axis} -- size of orbit (average radius)
    \item $e$: \textbf{Eccentricity} -- shape of orbit (0 = circle, $<$1 = ellipse); shown by ellipse elongation
    \item \textcolor{red!80!black}{$i$}: \textbf{Inclination} -- tilt angle between orbital plane and equatorial plane
    \item \textcolor{green!50!black}{$\Omega$}: \textbf{Right Ascension of Ascending Node (RAAN)} -- angle from $\gamma$ to ascending node
    \item \textcolor{brown!80!black}{$\omega$}: \textbf{Argument of periapsis} -- angle from ascending node to periapsis
    \item \textcolor{violet!80!black}{$\nu$}: \textbf{True anomaly} -- angle from periapsis to satellite (changes with time)
\end{itemize}
\end{tcolorbox}

\caption{Keplerian Orbital Elements in 3D perspective. The six elements ($a, e, i, \Omega, \omega, \nu$) uniquely define a satellite's position and velocity at any time.}
\label{fig:keplerian}
\end{figure}

\subsection{Two-Body Problem}

The fundamental equation of orbital motion:

\begin{equation}
\ddot{\mathbf{r}} = -\frac{\mu}{r^3} \mathbf{r}
\end{equation}

Where:
\begin{itemize}[leftmargin=*]
    \item $\mu = GM$ (gravitational parameter of Earth = 398,600 km$^3$/s$^2$)
    \item $\mathbf{r}$ = position vector
    \item $\ddot{\mathbf{r}}$ = acceleration
\end{itemize}

\begin{tcolorbox}[colback=yellow!10!white,colframe=emphcolor,title=Key Insight for Controls]
In orbit, there's no friction! Any velocity change is permanent unless you apply another force.
\end{tcolorbox}

\subsection{Hohmann Transfer}

To change orbits efficiently:
\begin{enumerate}
    \item Apply $\Delta v$ at periapsis to raise apoapsis
    \item Coast to new apoapsis
    \item Apply $\Delta v$ to circularize
\end{enumerate}

\textbf{Critical for your problem:} Rendezvous requires matching both position AND velocity!

\section{The Kessler Syndrome}

\subsection{The Cascading Collision Problem}

Named after NASA scientist Donald Kessler (1978), this is the critical threat:

\begin{enumerate}
    \item Two objects collide at $\sim$15 km/s relative velocity
    \item Collision creates thousands of new debris fragments
    \item Each fragment can cause more collisions
    \item \textbf{Exponential growth} of debris population
    \item Eventually, LEO becomes unusable (chain reaction)
\end{enumerate}

\subsection{Current Risk Level}

\begin{itemize}[leftmargin=*]
    \item We're approaching the critical density threshold
    \item Some orbits (800--1000 km) are already at risk
    \item Mega-constellations (Starlink: 42,000 satellites planned) increase risk
\end{itemize}

\subsection{Why Active Removal is Needed}

\begin{itemize}[leftmargin=*]
    \item Natural decay takes decades-to-centuries in LEO
    \item Atmospheric drag only significant below $\sim$600 km
    \item We need to remove $\sim$5--10 large objects per year to stabilize
\end{itemize}

\section{Spacecraft Dynamics \& Control in Space}

\subsection{Attitude Dynamics (Rotation)}

\textbf{Euler's Rotational Equation:}

\begin{equation}
\mathbf{I} \cdot \dot{\boldsymbol{\omega}} + \boldsymbol{\omega} \times (\mathbf{I} \cdot \boldsymbol{\omega}) = \boldsymbol{\tau}
\end{equation}

Where:
\begin{itemize}[leftmargin=*]
    \item $\mathbf{I}$ = Moment of inertia tensor (3$\times$3 matrix)
    \item $\boldsymbol{\omega}$ = Angular velocity vector
    \item $\boldsymbol{\tau}$ = Applied torque
\end{itemize}

\textbf{Attitude Representation:}
\begin{itemize}[leftmargin=*]
    \item Euler angles (intuitive but have singularities)
    \item Quaternions (4 parameters, no singularities -- preferred!)
    \item Rotation matrices (9 parameters, redundant)
\end{itemize}

\subsection{Actuators in Space}

\subsubsection{For Translation ($\Delta v$):}

\textbf{1. Thrusters}: Chemical or electric propulsion
\begin{itemize}[leftmargin=*]
    \item Impulse: $\Delta v = v_e \cdot \ln(m_0/m_f)$ -- Tsiolkovsky equation
    \item Fuel limited!
\end{itemize}

\subsubsection{For Rotation (torque):}

\begin{enumerate}
    \item \textbf{Reaction Wheels}: Spin internal mass, conserve angular momentum
    \begin{itemize}
        \item Fast, precise, can saturate (need desaturation)
    \end{itemize}

    \item \textbf{Control Moment Gyroscopes (CMGs)}: Gimbaled spinning wheels
    \begin{itemize}
        \item Higher torque than reaction wheels
    \end{itemize}

    \item \textbf{Magnetic Torquers}: Interact with Earth's magnetic field
    \begin{itemize}
        \item Slow, only certain directions, free propellant
    \end{itemize}

    \item \textbf{Thrusters}: Fast but uses fuel
\end{enumerate}

\subsection{Key Control Challenges}

\begin{itemize}[leftmargin=*]
    \item \textbf{Underactuated system}: Limited fuel, must be efficient
    \item \textbf{Coupled dynamics}: Translation-rotation coupling during capture
    \item \textbf{Long time delays}: Ground communication $\sim$1 second delay
    \item \textbf{Sensor noise}: Star trackers, gyros, GPS in space
\end{itemize}

\section{Rendezvous and Proximity Operations (RPO)}

\subsection{The Four Phases of Rendezvous}

\textbf{Phase 1: Phasing (Far Range, $>$100 km)}
\begin{itemize}[leftmargin=*]
    \item Match orbital plane (inclination)
    \item Adjust orbit to approach target
    \item Use ground-based tracking
\end{itemize}

\textbf{Phase 2: Approach (10--100 km)}
\begin{itemize}[leftmargin=*]
    \item Refine trajectory
    \item Begin autonomous sensing
    \item Switch to relative navigation
\end{itemize}

\textbf{Phase 3: Proximity Operations (100 m -- 10 km)}
\begin{itemize}[leftmargin=*]
    \item \textbf{Critical phase for your problem!}
    \item Careful trajectory planning
    \item Collision avoidance
    \item Continuous target tracking
\end{itemize}

\textbf{Phase 4: Close Proximity \& Capture ($<$100 m)}
\begin{itemize}[leftmargin=*]
    \item Final approach corridor
    \item Synchronize rotation rates
    \item Execute capture
\end{itemize}

\subsection{Hill-Clohessy-Wiltshire (HCW) Equations}

\textbf{Relative motion in orbit} (linearized):

\begin{align}
\ddot{x} - 2n\dot{y} - 3n^2x &= \frac{f_x}{m} \\
\ddot{y} + 2n\dot{x} &= \frac{f_y}{m} \\
\ddot{z} + n^2z &= \frac{f_z}{m}
\end{align}

Where:
\begin{itemize}[leftmargin=*]
    \item $(x, y, z)$: Local orbital frame (radial, along-track, cross-track)
    \item $n$: Mean motion (orbital angular velocity) = $\sqrt{\mu/a^3}$
    \item $\mathbf{f}$: Control forces
\end{itemize}

\begin{tcolorbox}[colback=blue!5!white,colframe=codeborder,title=Key Insight]
\begin{itemize}[leftmargin=*]
    \item Radial ($x$) and along-track ($y$) are \textbf{coupled} via Coriolis terms ($2n\dot{y}$, $2n\dot{x}$)
    \item Cross-track ($z$) is \textbf{decoupled}
    \item Natural motion creates elliptical relative trajectories
\end{itemize}
\end{tcolorbox}

\subsection{V-bar and R-bar Approaches}

\textbf{R-bar (Radial approach):}
\begin{itemize}[leftmargin=*]
    \item Approach from above/below
    \item Naturally stable against along-track drift
    \item Requires active station-keeping
\end{itemize}

\textbf{V-bar (Velocity vector approach):}
\begin{itemize}[leftmargin=*]
    \item Approach from behind/front
    \item Passive safety (natural separation if engines fail)
    \item Preferred for ISS
\end{itemize}

\section{Capture Mechanisms}

\subsection{Contact-Based Methods}

\subsubsection{1. Robotic Arm}
\begin{itemize}[leftmargin=*]
    \item Example: Canadarm2 on ISS, SSRMS
    \item \textbf{Pros}: Precise, proven technology
    \item \textbf{Cons}: Requires close proximity, complex dynamics
    \item \textbf{Control Challenge}: Free-floating base problem!
\end{itemize}

\begin{tcolorbox}[colback=yellow!10!white,colframe=emphcolor,title=Free-Floating Dynamics]
When arm moves $\rightarrow$ spacecraft base reacts (Newton's 3rd law)

You need \textbf{Generalized Jacobian} that accounts for base motion:
\begin{itemize}[leftmargin=*]
    \item Position of end-effector depends on both joint angles AND base motion
    \item Captures induce momentum transfer $\rightarrow$ spacecraft tumbles
\end{itemize}
\end{tcolorbox}

\subsubsection{2. Harpoon}
\begin{itemize}[leftmargin=*]
    \item Fire penetrating spike into target
    \item Tether connects chaser to target
    \item \textbf{Pros}: Works for non-cooperative targets
    \item \textbf{Cons}: Can fragment target, tether dynamics complex
\end{itemize}

\subsubsection{3. Net}
\begin{itemize}[leftmargin=*]
    \item Deploy net to envelop target
    \item \textbf{Pros}: Gentle, encapsulates fragments
    \item \textbf{Cons}: Deployment dynamics, tangling risk
\end{itemize}

\subsubsection{4. Gripper/Grapple}
\begin{itemize}[leftmargin=*]
    \item Mechanical jaws grip target
    \item Requires attachment point (cooperative targets)
\end{itemize}

\subsection{Contactless Methods}

\subsubsection{1. Electrostatic Tractor}
\begin{itemize}[leftmargin=*]
    \item Charge spacecraft oppositely $\rightarrow$ Coulomb attraction
    \item \textbf{Pros}: No contact, scalable force
    \item \textbf{Cons}: Requires charge management, works at close range
\end{itemize}

\subsubsection{2. Laser Ablation}
\begin{itemize}[leftmargin=*]
    \item Vaporize surface material $\rightarrow$ creates thrust on target
    \item \textbf{Pros}: Standoff distance, no rendezvous needed
    \item \textbf{Cons}: Debris from ablation, power requirements
\end{itemize}

\subsubsection{3. Ion Beam Shepherd}
\begin{itemize}[leftmargin=*]
    \item Plasma thruster creates force on target
    \item \textbf{Pros}: Gentle, continuous force
    \item \textbf{Cons}: Very slow, complex plasma dynamics
\end{itemize}

\subsubsection{4. Magnetic Tractor}
\begin{itemize}[leftmargin=*]
    \item Use magnetic fields (for ferromagnetic debris)
    \item \textbf{Pros}: No contact
    \item \textbf{Cons}: Only works for magnetic materials
\end{itemize}

\section{Control Challenges for Non-Cooperative Targets}

\subsection{What Makes a Target ``Non-Cooperative''?}

\begin{enumerate}
    \item \textbf{Unknown dynamics}: Tumbling rate unknown
    \item \textbf{No communication}: Can't command it
    \item \textbf{Unknown mass properties}: Center of mass, inertia tensor unknown
    \item \textbf{Unknown shape/structure}: Visual only
    \item \textbf{Unpredictable}: Solar radiation pressure, atmospheric drag variations
\end{enumerate}

\subsection{State Estimation Problem}

You need to estimate:
\begin{itemize}[leftmargin=*]
    \item \textbf{Position \& velocity} (6 DOF translation)
    \item \textbf{Orientation \& angular velocity} (6 DOF rotation)
    \item \textbf{Total: 12-state estimation problem}
\end{itemize}

\textbf{Sensors:}
\begin{itemize}[leftmargin=*]
    \item \textbf{LIDAR}: Range and range-rate
    \item \textbf{Cameras}: Visual tracking, structure from motion
    \item \textbf{Radar}: All-weather tracking
\end{itemize}

\textbf{Estimation Algorithms:}
\begin{itemize}[leftmargin=*]
    \item Extended Kalman Filter (EKF)
    \item Unscented Kalman Filter (UKF) -- better for nonlinear
    \item Particle Filters -- for highly nonlinear/multimodal
\end{itemize}

\subsection{Tumbling Dynamics}

Debris often tumbles due to:
\begin{itemize}[leftmargin=*]
    \item Collision-induced rotation
    \item Gravity gradient torque
    \item Solar radiation pressure
    \item Magnetic field interaction
\end{itemize}

\textbf{Euler's Equation} (torque-free motion):

\begin{align}
I_1\dot{\omega}_1 &= (I_2 - I_3)\omega_2\omega_3 \\
I_2\dot{\omega}_2 &= (I_3 - I_1)\omega_3\omega_1 \\
I_3\dot{\omega}_3 &= (I_1 - I_2)\omega_1\omega_2
\end{align}

\textbf{Result:} Complex 3D tumbling, periodic for rigid bodies

\subsection{Detumbling Strategies}

\textbf{Option 1: Match Rotation}
\begin{itemize}[leftmargin=*]
    \item Estimate target angular velocity
    \item Synchronize chaser rotation
    \item Capture in rotating frame
    \item \textbf{Challenge}: High-rate tumbling $\rightarrow$ high control effort
\end{itemize}

\textbf{Option 2: Active Detumbling Before Capture}
\begin{itemize}[leftmargin=*]
    \item Apply external torque (contactless methods)
    \item Slow down tumbling gradually
    \item Then capture
    \item \textbf{Challenge}: Takes time, requires precise torque control
\end{itemize}

\textbf{Option 3: Capture While Tumbling}
\begin{itemize}[leftmargin=*]
    \item Fast capture mechanism
    \item Accept the momentum transfer
    \item Detumble combined system afterward
    \item \textbf{Challenge}: Huge impact forces, risk of damage
\end{itemize}

\section{Deorbiting Techniques}

\subsection{Natural Decay}

\begin{itemize}[leftmargin=*]
    \item Atmospheric drag gradually lowers orbit
    \item Decay time = $f$(altitude, area-to-mass ratio, solar activity)
    \item \textbf{Problem:} Can take decades above 600 km
\end{itemize}

\subsection{Active Deorbiting}

\subsubsection{1. Direct Re-entry}
\begin{itemize}[leftmargin=*]
    \item Use propulsion to lower periapsis below $\sim$80 km
    \item Atmospheric drag increases exponentially
    \item Object burns up
    \item \textbf{$\Delta v$ required:} $\sim$100--200 m/s from LEO
\end{itemize}

\subsubsection{2. Drag Augmentation}

\textbf{a) Drag Sail}
\begin{itemize}[leftmargin=*]
    \item Deploy large surface area (sail)
    \item Increases drag force: $F_{\text{drag}} = 0.5 \rho v^2 C_D A$
    \item Accelerates natural decay
    \item \textbf{Pros:} Passive, no propellant
    \item \textbf{Cons:} Attitude stabilization needed
\end{itemize}

\textbf{b) Inflatable Structures}
\begin{itemize}[leftmargin=*]
    \item Similar to drag sail but lighter
    \item Can increase area by 10--100$\times$
\end{itemize}

\textbf{c) Electrodynamic Tether}
\begin{itemize}[leftmargin=*]
    \item Long conducting wire in Earth's magnetic field
    \item Generates current $\rightarrow$ Lorentz force $\rightarrow$ drag
    \item \textbf{Pros:} No propellant, can generate power
    \item \textbf{Cons:} Complex deployment, tether dynamics
\end{itemize}

\subsubsection{3. Graveyard Orbit}
\begin{itemize}[leftmargin=*]
    \item Boost to higher orbit ($>$2000 km)
    \item Out of useful LEO region
    \item \textbf{Problem:} Doesn't solve long-term issue
\end{itemize}

\subsection{Controlled Re-entry}

For large objects ($>$1 ton):
\begin{itemize}[leftmargin=*]
    \item Target specific ocean areas (South Pacific)
    \item Prevent casualty on ground
    \item Requires precise trajectory control
\end{itemize}

\section{Key Differences from Terrestrial Robotics}

\begin{enumerate}
    \item \textbf{No Fixed Base}
    \begin{itemize}
        \item Manipulator motion $\rightarrow$ base reaction
        \item Need \textbf{Generalized Jacobian} and \textbf{Dynamic Singularities}
        \item Momentum conservation critical
    \end{itemize}

    \item \textbf{No Gravity (Microgravity)}
    \begin{itemize}
        \item Objects don't ``fall''
        \item Small forces persist
        \item Contact forces dominate
    \end{itemize}

    \item \textbf{Orbital Dynamics}
    \begin{itemize}
        \item Reference frame accelerates (centripetal)
        \item Coriolis effects in relative motion
        \item Tidal forces (gravity gradient)
    \end{itemize}

    \item \textbf{Propellant Constraints}
    \begin{itemize}
        \item Every maneuver costs fuel
        \item Fuel mass = mission lifetime
        \item Optimization critical
    \end{itemize}

    \item \textbf{Extreme Environment}
    \begin{itemize}
        \item Temperature extremes: $-150^\circ$C to $+150^\circ$C
        \item Radiation damages electronics
        \item Vacuum $\rightarrow$ no convective cooling
        \item Atomic oxygen erodes materials
    \end{itemize}

    \item \textbf{Communication Delays}
    \begin{itemize}
        \item Ground stations: $\sim$1 sec delay
        \item Requires autonomy
        \item Can't teleoperate fast maneuvers
    \end{itemize}

    \item \textbf{Sensing Challenges}
    \begin{itemize}
        \item GPS limited to LEO
        \item Visual tracking affected by lighting (sun/eclipse)
        \item No tactile feedback before contact
    \end{itemize}
\end{enumerate}

\section{Control Architecture for ADR Mission}

\subsection{Typical System Layers}

\begin{center}
\begin{minipage}{0.8\textwidth}
\begin{tcolorbox}[colback=blue!5!white,colframe=codeborder]
\centering
\textbf{Mission Planning (Ground/Onboard)} \\
Target selection, Trajectory optimization

\vspace{0.3cm}
$\downarrow$

\vspace{0.3cm}
\textbf{Guidance (Path Planning)} \\
Collision avoidance, Approach corridors, Fuel optimization

\vspace{0.3cm}
$\downarrow$

\vspace{0.3cm}
\textbf{Navigation (State Estimation)} \\
Sensor fusion (LIDAR, camera), Target pose estimation, \\
Relative state (position, velocity)

\vspace{0.3cm}
$\downarrow$

\vspace{0.3cm}
\textbf{Control (Low-Level)} \\
Attitude control (PD, MPC), Thruster firing logic, \\
Manipulator control
\end{tcolorbox}
\end{minipage}
\end{center}

\subsection{Control Algorithms Relevant for This Problem}

\textbf{1. PID Control}
\begin{itemize}[leftmargin=*]
    \item Simple, robust
    \item Works for attitude stabilization
    \item May struggle with coupled dynamics
\end{itemize}

\textbf{2. Model Predictive Control (MPC)}
\begin{itemize}[leftmargin=*]
    \item Optimal for constrained systems
    \item Handles thruster saturation
    \item Fuel optimization
    \item Collision avoidance constraints
    \item \textbf{Highly recommended for proximity ops!}
\end{itemize}

\textbf{3. Sliding Mode Control}
\begin{itemize}[leftmargin=*]
    \item Robust to uncertainties
    \item Good for non-cooperative targets
    \item Chattering can waste fuel
\end{itemize}

\textbf{4. Adaptive Control}
\begin{itemize}[leftmargin=*]
    \item Estimate unknown target parameters online
    \item Adjust control gains
    \item Useful for tumbling targets
\end{itemize}

\textbf{5. Optimal Control (LQR/LQG)}
\begin{itemize}[leftmargin=*]
    \item Minimize fuel + time
    \item Well-suited for rendezvous
\end{itemize}

\section{Recommended Focus Areas for Your Solution}

\subsection{As a Robotics/Controls Engineer, You Can Contribute:}

\begin{enumerate}
    \item \textbf{Robust State Estimation}
    \begin{itemize}
        \item Develop filters for tumbling target tracking
        \item Sensor fusion (vision + LIDAR)
        \item Handle occlusions, lighting changes
    \end{itemize}

    \item \textbf{Trajectory Planning}
    \begin{itemize}
        \item Fuel-optimal approach trajectories
        \item Collision avoidance
        \item Reachability analysis
    \end{itemize}

    \item \textbf{Capture Control}
    \begin{itemize}
        \item Free-floating base control
        \item Impact minimization
        \item Post-capture stabilization
    \end{itemize}

    \item \textbf{Detumbling Strategies}
    \begin{itemize}
        \item Contactless detumbling (if using ion beam, etc.)
        \item Momentum management
        \item Combined system control after capture
    \end{itemize}

    \item \textbf{Fault Tolerance}
    \begin{itemize}
        \item Thruster failures
        \item Sensor degradation
        \item Safe-mode behaviors
    \end{itemize}
\end{enumerate}

\section{Key Equations Cheat Sheet}

\subsection{Orbital Mechanics}

\begin{align}
\text{Orbital velocity:} \quad v &= \sqrt{\frac{\mu}{r}} \\
\text{Orbital period:} \quad T &= 2\pi\sqrt{\frac{a^3}{\mu}} \\
\text{Vis-viva equation:} \quad v^2 &= \mu\left(\frac{2}{r} - \frac{1}{a}\right) \\
\text{Hohmann } \Delta v: \quad \Delta v_1 &= \sqrt{\frac{\mu}{r_1}}\left(\sqrt{\frac{2r_2}{r_1+r_2}} - 1\right)
\end{align}

\subsection{Relative Motion (HCW)}

\begin{align}
\ddot{x} - 2n\dot{y} - 3n^2x &= \frac{f_x}{m} \\
\ddot{y} + 2n\dot{x} &= \frac{f_y}{m} \\
\ddot{z} + n^2z &= \frac{f_z}{m}
\end{align}

where $n = \sqrt{\mu/a^3}$

\subsection{Attitude Dynamics}

\begin{align}
\mathbf{I} \cdot \dot{\boldsymbol{\omega}} + \boldsymbol{\omega} \times (\mathbf{I} \cdot \boldsymbol{\omega}) &= \boldsymbol{\tau} \\
\text{Quaternion kinematics:} \quad \dot{\mathbf{q}} &= 0.5 \, \Omega(\boldsymbol{\omega})\mathbf{q}
\end{align}

\subsection{Drag Force}

\begin{equation}
F_{\text{drag}} = 0.5 \rho v^2 C_D A
\end{equation}

where $\rho$ = atmospheric density

\subsection{Tsiolkovsky Rocket Equation}

\begin{equation}
\Delta v = v_e \ln\left(\frac{m_0}{m_f}\right)
\end{equation}

where $v_e$ = exhaust velocity

\section{Resources for Further Learning}

\subsection{Papers (from problem statement)}

\begin{enumerate}
    \item A. Ledkov, V. Aslanov, ``Review of contact and contactless active space debris removal approaches'', \textit{Prog. Aerosp. Sci.} 134 (2022) 100858
    \item M. Shan, J. Guo, E. Gill, ``Review and comparison of active space debris capture and removal methods'', \textit{Prog. Aero. Sci.} 80 (2016) 18-32
\end{enumerate}

\subsection{Recommended Textbooks}

\begin{itemize}[leftmargin=*]
    \item \textbf{Orbital Mechanics}: ``Orbital Mechanics for Engineering Students'' by Howard Curtis
    \item \textbf{Spacecraft Dynamics}: ``Spacecraft Dynamics and Control'' by Marcel Sidi
    \item \textbf{Space Robotics}: ``Robotics and Automation in Space'' -- Yoshida/Umetani
\end{itemize}

\subsection{Online Resources}

\begin{itemize}[leftmargin=*]
    \item ESA Space Debris Office: \url{https://www.esa.int/Space_Safety/Space_Debris}
    \item NASA Orbital Debris Program: \url{https://orbitaldebris.jsc.nasa.gov/}
\end{itemize}

\subsection{Software Tools}

\begin{itemize}[leftmargin=*]
    \item \textbf{GMAT} (General Mission Analysis Tool): Free trajectory design
    \item \textbf{STK} (Systems Tool Kit): Commercial, orbit visualization
    \item \textbf{Orekit}: Open-source orbital mechanics library (Python/Java)
\end{itemize}

\section{Summary: Your Robotics Skills Applied to Space}

\begin{table}[h]
\centering
\begin{tabular}{|p{6cm}|p{6cm}|}
\hline
\textbf{Terrestrial Robotics} & \textbf{Space Robotics (ADR)} \\
\hline
Fixed base manipulator & Free-floating base \\
\hline
Cooperative targets & Non-cooperative, tumbling \\
\hline
Gravity assists & Orbital dynamics dominate \\
\hline
Unlimited power & Solar/battery constrained \\
\hline
Rich sensing & Limited, delayed sensing \\
\hline
Ground-based testing & Simulation + limited space testing \\
\hline
Reaction forces to ground & Momentum conservation \\
\hline
\end{tabular}
\end{table}

\vspace{1cm}

\begin{tcolorbox}[colback=green!5!white,colframe=green!50!black,title=Your Advantage]
\textbf{Control theory, state estimation, path planning, and optimization translate directly!} The physics is different but the mathematical frameworks are similar.
\end{tcolorbox}

\vspace{1cm}

\begin{center}
{\Large \textbf{Good luck with your hackathon!}}
\end{center}

\end{document}
